\documentclass[aspectratio=169, 12pt]{beamer}

% --- Language and Fonts ---
\usepackage{fontspec}
\usepackage[english]{babel}
\usepackage{amsmath, amssymb, bm}
% tikz
\usepackage{tikz}
\usetikzlibrary{shapes, arrows.meta, positioning, shadows, calc}

% --- Bibliography (BibTeX + natbib) ---
\usepackage{natbib}
\bibliographystyle{chicago} 

% =========================================
% ===== keybox ====================
% =========================================
\usepackage[most]{tcolorbox}
\newtcolorbox{keybox}[1]{
  enhanced,
  colback=headbg!4,
  colframe=headbg,
  coltitle=headfg,
  colbacktitle=headbg,
  title=#1,
  fonttitle=\bfseries,
  boxrule=0.8pt,
  arc=1.5mm,
  left=2mm,
  right=2mm,
  top=1mm,
  bottom=1mm,
  before skip=0.2cm,
  after skip=0.2cm
}

\usepackage{booktabs}

% =========================================
% ===== Custom Footline (Page Number) =====
% =========================================


\setbeamercolor{footline}{fg=headbg}
\setbeamerfont{footline}{size=\scriptsize}

\setbeamertemplate{footline}{%
  \hfill%
  \usebeamercolor[fg]{footline}%
  \usebeamerfont{footline}%
  \insertframenumber\,/\,\inserttotalframenumber\kern1em\vskip3pt%
}

\setbeamertemplate{navigation symbols}{}

\usetheme{default}
\usepackage{xcolor}

\definecolor{headbg}{RGB}{30, 60, 120}
\definecolor{headfg}{RGB}{255, 255, 255}
\definecolor{accent}{RGB}{200, 50, 50}

\setbeamercolor{headline}{bg=headbg, fg=headfg}
\setbeamercolor{section in head/foot}{bg=headbg, fg=headfg}
\setbeamercolor{subsection in head/foot}{bg=headbg!80, fg=headfg}
\setbeamercolor{frametitle}{bg=headbg, fg=headfg}

% Step indicator command
\newcommand{\stepindicator}[1]{%
  \hfill {\scriptsize \textit{\textcolor{headbg}{Step #1 / 6}}}%
}
\DeclareMathOperator*{\argmin}{arg\,min}

\usepackage{hyperref}


% --- Title ---
\title[]{A First Course in Causal Inference Ch.4}
\author[Kodai Minesaki]{Kodai Minesaki}
\institute[Faculty of Engineering, The University of Tokyo]{Faculty of Engineering, The University of Tokyo}
\date{\today}

\begin{document}

% --- Title Slide ---
\begin{frame}
\titlepage
\end{frame}

\begin{frame}
  
  {
    \large
    
    Ch.4 Neymanian Repeated Sampling Inference in CRE
  }

\end{frame}

\begin{frame}{Ch.4 Neymanian Repeated Sampling Inference in CRE}
  \begin{itemize}
    \item \citet{neyman1923application} has proposed an unbiased point estimator and a conservative confidence interval based on the sampling distribution of the point estimator
    \item This chapter covers \citet{neyman1923application}'s fundamental parts.
  \end{itemize}
\end{frame}

\begin{frame}{4.1 Finite population quantities}
  \begin{itemize}
    \item Setup: CRE with $n$ units (treatement group: $n_1$ units, control group: $n_0$ units); For $i=1,...,n$, potential outcomes $Y_i(1)$ and $Y_i(0)$ and individual effect $\tau_i = Y_i(1)-Y_i(0)$.
    \item We then obtain finite population means:
    \[
    \overline{Y}(1) = \frac{1}{n}\sum_{i=1}^{n}Y_i(1), \quad \overline{Y}(0) = \frac{1}{n}\sum_{i=1}^{n}Y_i(0),
    \]
    \item variances:
    \[
    S^2(1) = \frac{1}{n-1}\sum_{i=1}^{n}\{Y_i(1)-\overline{Y}(1)\}^2, \quad S^2(0) = \frac{1}{n-1}\sum_{i=1}^{n}\{Y_i(0)-\overline{Y}(0)\}^2,
    \]
    \item and covariance:
    \[
    S(1,0) = \frac{1}{n-1}\sum_{i=1}^{n}\{Y_i(1)-\overline{Y}(1)\}\{Y_i(0)-\overline{Y}(0)\}.
    \]
  \end{itemize}
\end{frame}

\begin{frame}{4.1 Finite population quantities}
  \begin{itemize}
    \item The individual effects have mean
    \[
    \tau = \frac{1}{n}\sum_{i=1}^{n}\tau_i = \overline{Y}(1)-\overline{Y}(0), 
    \]
    \item and variance
    \[
    S^2(\tau) = \frac{1}{n-1}\sum_{i=1}^{n}(\tau_i-\tau)^2
    \]
  \end{itemize}
  We then have
  \begin{keybox}{Lemma 4.1}
    \[
    2S(1,0) = S^2(1)+S^2(0) - S^2(\tau).
    \]
  \end{keybox}
\end{frame}

\begin{frame}{4.2 \citet{neyman1923application}'s theorem}
  \begin{itemize}
    \item We are interested in estimating the average causal effect $\tau$ based on the (observed) data $(Z_i, Y_i)_{i=1}^{n}$ from a CRE.
    \item Based on the observed data, the sample means:
    \[
    \widehat{\overline{Y}}(1) = \frac{1}{n_1}\sum_{i=1}^{n}Z_iY_i, \quad \widehat{\overline{Y}}(0) = \frac{1}{n_0}\sum_{i=1}^{n}(1-Z_i)Y_i,
    \]
    \item and the sample variances:
    \[
    \widehat{S}^2(1) = \frac{1}{n_1-1}\sum_{i=1}^{n}Z_i\{Y_i-\widehat{\overline{Y}}(1)\}^2, \quad \widehat{S}^2(0) = \frac{1}{n_0-1}\sum_{i=1}^{n}(1-Z_i)\{Y_i-\widehat{\overline{Y}}(0)\}^2.
    \]
    \item Obviously, there are no sample covariance and $S^2(\tau)$ because \textbf{we cannot observe both potential outcomes for each unit}.
  \end{itemize}
\end{frame}

\begin{frame}{4.2 \citet{neyman1923application}'s theorem}
  \begin{keybox}{Theorem 4.1}
    \begin{enumerate}
      \item the difference-in-means estimator $\widehat{\tau}=\widehat{\overline{Y}}(1)-\widehat{\overline{Y}}(0)$ is unbiased for $\tau$:$$\mathbb{E}[\widehat{\tau}]=\tau.$$
      \item $\widehat{\tau}$ has variance
      \[
      \operatorname{var}(\widehat{\tau}) = \frac{S^2(1)}{n_1} + \frac{S^2(0)}{n_0} - \frac{S^2(\tau)}{n} = \frac{n_0}{n_1n}S^2(1)+\frac{n_1}{n_0n}S^2(0) + \frac{2}{n}S(1,0).
      \]
      \item the variance estimator $\widehat{V} = \frac{\widehat{S}^2(1)}{n_1} + \frac{\widehat{S}^2(0)}{n_0}$ is conservative for estimating $\operatorname{var}(\widehat{\tau})$ in the sense that $\mathbb{E}[\widehat{V}]-\operatorname{var}(\widehat{\tau}=\frac{S^2(\tau)}{n})\geq0$ with equality holding if and only if $\tau_i=\tau$ for all units.
    \end{enumerate}
  \end{keybox}
\end{frame}

\begin{frame}{4.2 \citet{neyman1923application}'s theorem}
  \begin{itemize}
    \item The FRT works for any test statistic but \citet{neyman1923application}'s theorem is only about the difference in means.
    \item If $S(1,0)>0$, then those treated units also have relatively large control potential outcomes\footnote{In other words, large control potential outcomes are not ``observed.''}. $\leadsto$ the observed outcomes under control are relatively small $\leadsto$ large $\widehat{\tau}$.
    \item If $S(1,0)<0$, the those treated units have relatively small control potential outcomes\footnote{In other words, small control potential outcomes are not ``observed.''}. $\leadsto$ the observed outcomes under control are relatively large $\leadsto$ small $\widehat{\tau}$.
  \end{itemize}
\end{frame}

\begin{frame}{4.2 \citet{neyman1923application}'s theorem}
  \begin{keybox}{Thoerem 4.2}
    Let $n\to \infty$ and $n_i\to \infty$. If $\dfrac{n_1}{n}$ has a limiting value in $(0,1)$, $\{S^2(1), S^2(0), S(1,0)\}$ have limiting values, and 
    \begin{equation}
      \max_{1\leq i\leq n}\frac{\{Y_i(1)-\overline{Y}(1)\}^2}{n}\to 0,\quad \max_{1\leq i\leq n}\frac{\{Y_i(1)-\overline{Y}(0)\}^2}{n}\to 0,
      \label{eq:thm4.2-1}
    \end{equation}
    then 
    \begin{equation}
      \frac{\widehat{\tau}-\tau}{\sqrt{\operatorname{var}(\widehat{\tau})}}\xrightarrow[d]{}\mathcal{N}(0,1),
      \label{eq:thm4.2-2}
    \end{equation}
    and
    \begin{equation}
      \widehat{S}^2(1)\xrightarrow[p]{}S^2(1), \quad \widehat{S}^2(1)\xrightarrow[p]{}S^2(0)
      \label{eq:thm4.2-3}
    \end{equation}
    \label{thm4.2}
  \end{keybox}
\end{frame}

\begin{frame}{4.2 \citet{neyman1923application}'s theorem}
  The intuitive understanding for Thm.~4.2 is:
  \begin{itemize}
    \item Eq.~\eqref{eq:thm4.2-1}: potential outcomes for each unit are not extremely large to dominate the variability\footnote{Formally, $\max_{1\leq i \leq n}\|Y_i(1)-\overline{Y}(1)\|=o(\sqrt{n})$.}. 
    \item Suppose a case one unit has an extremely large outcome, then $\widehat{\tau}$ depends on which group the unit belongs to: Eq.~\eqref{eq:thm4.2-2} does not hold.
    \item Eq.~\eqref{eq:thm4.2-3} ensures \citet{neyman1923application}'s conservative CI for $\tau$, because $\widehat{V}\geq \operatorname{var}(\widehat{\tau})$.
  \end{itemize}
\end{frame}

\begin{frame}{4.3 Proofs}
  Unbiasedness of $\widehat{\tau}$:
  {\small
  \begin{align*}
    \widehat{\tau} 
    = \frac{1}{n_1}\sum_{i=1}^{n}Z_iY_i - \frac{1}{n_0}\sum_{i=1}^{n}(1-Z_i)Y_i
    =\frac{1}{n_1}\sum_{i=1}^{n}Z_iY_i(1) - \frac{1}{n_0}\sum_{i=1}^{n}(1-Z_i)Y_i(0),
  \end{align*}
  then
  \begin{align*}
    \mathbb{E}[\widehat{\tau}] 
    &=\mathbb{E}\left[\frac{1}{n_1}\sum_{i=1}^{n}Z_iY_i(1) - \frac{1}{n_0}\sum_{i=1}^{n}(1-Z_i)Y_i(0)\right] \\
    &=\frac{1}{n_1}\sum_{i=1}^{n}\mathbb{E}[Z_i]\,Y_i(1) - \frac{1}{n_0}\sum_{i=1}^{n}\mathbb{E}[1-Z_i]\,Y_i(0) \\
    &= \frac{1}{n_1}\sum_{i=1}^{n}\frac{n_1}{n}Y_i(1) - \frac{1}{n_0}\sum_{i=1}^{n}\frac{n_0}{n}Y_i(0)
    = \frac{1}{n}\sum_{i=1}^{n}Y_i(1) - \frac{1}{n}\sum_{i=1}^{n}Y_i(0) = \tau. \quad \Box
  \end{align*}
  }
\end{frame}

\begin{frame}{4.3 Proofs}
  $\operatorname{var}(\widehat{\tau})$:
  {\small
  We write $\widehat{\tau}$ as
  \[
  \widehat{\tau} = \sum_{i=1}^{n}Z_i\left\{\frac{Y_i(1)}{n_1}+\frac{Y_i(0)}{n_0}\right\} - \frac{1}{n_0}\sum_{i=1}^{n}Y_i(0).
  \]
  Then, the variance of $\widehat{\tau}$ follows:
  \begin{align*}
    \operatorname{var}(\widehat{\tau}) 
    &= \frac{n_1n_0}{n(n-1)}\sum_{i=1}^{n}\left\{\frac{Y_i(1)}{n_1}+\frac{Y_i(0)}{n_0}-\frac{\overline{Y}(1)}{n_1}-\frac{\overline{Y}(0)}{n_0}\right\}^2\\
    &= \frac{n_1n_0}{n(n-1)}\left[\frac{1}{n_1^2}\sum_{i=1}^{n}\{Y_i(1)-\overline{Y}(1)\}^2+\frac{1}{n_0^2}\sum_{i=1}^{n}\{Y_i(0)-\overline{Y}(0)\}^2\right.\\
    &\qquad\left.+\frac{2}{n_1n_0}\sum_{i=1}^{n}\{Y_i(1)-\overline{Y}(1)\}\{Y_i(0)-\overline{Y}(0)\}\right]\\
    &= \frac{n_0}{n_1n}S^2(1) + \frac{n_1}{n_0n}S^2(0) + \frac{2}{n}S(1,0). 
  \end{align*}
  }
\end{frame}

\begin{frame}{4.3 Proofs}

    Here, the first equality holds because of the following Lemma:
    \begin{keybox}{Lemma C.2}
      Under simple random sampling, the sample means are unbiased for the population means:
      \[
      \mathbb{E}\left[\widehat{\overline{c}}\right] = \overline{c}, \quad \mathbb{E}\left[\widehat{\overline{d}}\right] = \overline{d}.
      \]
      Their variances and covariance are
      \[
      \operatorname{var}\left(\widehat{\overline{c}}\right) = \frac{n_0}{nn_1}S_c^2,\quad \operatorname{var}\left(\widehat{\overline{d}}\right) = \frac{n_0}{nn_1}S_d^2, \quad \operatorname{cov}\left(\widehat{\overline{c}}, \widehat{\overline{d}}\right) = \frac{n_0}{nn_1}S_cd.
      \]
    \end{keybox}

\end{frame}

\begin{frame}{4.3 Proofs}
  Let $A_i = \dfrac{Y_i(1)}{n_1}+\dfrac{Y_i(0)}{n_0}$, then $\operatorname{var}(\widehat{\tau}) = \operatorname{var}\left(\sum_{i=1}^{n}Z_iA_i\right)$. In a CRE, we draw a simple random sample of size $n_1$ for a treatment group. Then, $\overline{A} = \dfrac{1}{n_1}\sum_{_i=1}^{n}Z_iA_i.$ Following Lemma~C.2, $\operatorname{var}(\overline{A})=\dfrac{n_0}{nn_1}S_A^2$. Hence, 
  \begin{align*}
    \operatorname{var}(\widehat{\tau}) = \operatorname{var}\left(\sum_{i=1}^{n}Z_iA_i\right) 
    = \operatorname{var}(n_1\overline{A}) = n_1^2\operatorname{\overline{A}}
    = n_1^2 \frac{n_0}{nn_1}S_A^2 = \frac{n_0n_1}{n}S_A^2.
  \end{align*}
  Because $S_A^2 = \dfrac{1}{n-1}\sum_{i=1}^{n}\{A_i-\overline{A}\}^2$, we obtain $\operatorname{var}(\widehat{\tau}) = \dfrac{n_0n_1}{n(n-1)}\sum_{i=1}^{n}\left\{\dfrac{Y_i(1)}{n_1}+\dfrac{Y_i(0)}{n_0}-\dfrac{\overline{Y}(1)}{n_1}-\dfrac{\overline{Y}(0)}{n_0}\right\}^2$
\end{frame}

\begin{frame}{4.3 Proofs}
  From Lemma~4.1, we also write the variance as 
  \[
  \operatorname{var}(\widehat{\tau}) = \frac{n_0}{n_1n}S^2(1) + \frac{n_1}{n_0n}S^2(0) + \frac{1}{n}(S^2(1)+S^2(0)-S^2(\tau)) = \frac{S^2(1)}{n_1} + \frac{S^2(0)}{n_0} - \frac{S^2(\tau)}{n}. \quad \Box
  \]
\end{frame}

\begin{frame}{4.4 Regression analysis of the CRE}
  \begin{itemize}
    \item Practitioners often use regression-based inference for the average causal effect $\tau$.
    \item A standard approach is OLS: $$(\widehat{\alpha}, \widehat{\beta}) = \argmin_{(a,b)}\sum_{i=1}^{n}(Y_i-a-bZ_i)^2.$$
    \item Here, $\widehat{\tau} = \widehat{\beta}$.
    \item However, the usual variance estimator from the OLS equals
    \[
    \widehat{V}_\text{OLS} = \frac{n(n_1-1)}{(n-2)n_1n_0}\widehat{S}^2(1) + \frac{n(n_0-1)}{(n-2)n_1n_0}\widehat{S}^2(0) \approx \frac{\widehat{S}^2(1)}{n_0} + \frac{\widehat{S}^2(0)}{n_1},
    \]
    where the approximation holds with large $n_1$ and $n_0$ (it can differ).
    \item EHW robust variance estimator is closer to $\widehat{V}$ and HC2 varianct of the EHW robust variance estimator is identical to $\widehat{V}$.
  \end{itemize}
\end{frame}



% =========================

% =========================
\section{}
\begin{frame}[allowframebreaks, noframenumbering]{References}
  \beamerdefaultoverlayspecification{}
  \small
  \bibliography{../references}
\end{frame}

\end{document}